\subsection*{Passes 4324--4334: the chamber Hecke machine and an exact frontier correction}

\paragraph{4324: two native chamber switches.}
On the $160$ incident point--line chambers, let $P$ change the line while holding the
point fixed and let $L$ change the point while holding the line fixed.  Each is forty
disjoint copies of the $K_4$ adjacency operator, and GAP proves
\[
 P^2=2P+3I,\qquad L^2=2L+3I,\qquad PLPL=LPLP.
\]
The eight standard type-$C_2$ words are linearly independent.  Thus this chamber carrier
realizes the full eight-dimensional $q=3$ Iwahori--Hecke image.  Its exact module ledger is
\[
 1\chi_{3,3}+15\chi_{3,-1}+15\chi_{-1,3}+81\chi_{-1,-1}+24V_2,
\]
which explains the long-standing flag-scheme dimensions $1+30+48+81=160$.

\paragraph{4325--4326: the Hashimoto turn and its chirality.}
With the $320$ oriented Levi edges ordered by orientation,
\[
 B_{\rm Levi}=\begin{pmatrix}0&L\\P&0\end{pmatrix},\qquad
 B_{\rm Levi}^2=\operatorname{diag}(LP,PL).
\]
The two products are disjoint binary $9$-out halves of the chamber distance-two shell.
For $K=LP$,
\[
 \chi_K(x)=(x-9)(x-1)^{81}(x+3)^{30}(x^2+9)^{24}.
\]
The antisymmetric turn $\Omega=LP-PL$ has
\[
 \chi_\Omega(x)=x^{112}(x^2+60)^{24},\quad
 \Pi_{48}=-\Omega^2/60,\quad \operatorname{rank}\Pi_{48}=48.
\]
Consequently $\Omega/\sqrt{60}$ is a canonical complex structure on the conjugate
$24+24$ chamber channel.  This is a finite operator theorem, not a $W(G_2)$ action.

\paragraph{4327: exact folded-propagator normal form.}
Let $C=P+L$, $X=(C-2I)\Pi_{48}$, and let $F$ be the old folded cubic point-Hashimoto
operator compressed to $\Pi_{48}$.  In the four-dimensional packet algebra
$\{\Pi_{48},X,\Omega,X\Omega\}$,
\[
 F=-68\Pi_{48}-31X-\frac{21}{2}\Omega+\frac23X\Omega,
 \qquad (F+68\Pi_{48})^2=-689\Pi_{48}.
\]
The old numerical $-6455$ is now explained exactly by the square of the off-diagonal
part.  No continuum propagator, particle, mass, or coupling identification follows.

\paragraph{4334: the $24+24$ count becomes an object-level carrier split.}
Let $E_p,E_\ell$ be the eigenvalue-$2$ projectors of the W33 point graph and its
dual line graph, and pull them to the chamber space as orthogonal rank-$24$ projectors
$Q_p,Q_\ell$.  GAP proves
\[
 \operatorname{im}\Pi_{48}
 =\operatorname{im}Q_p\oplus\operatorname{im}Q_\ell,\qquad
 Q_pQ_\ell Q_p=\frac38Q_p,\quad
 Q_\ell Q_pQ_\ell=\frac38Q_\ell.
\]
Thus the two natural summands meet only in zero but are uniformly nonorthogonal: all
$24$ principal angles have cosine $\sqrt6/4$.  With $Q=Q_p+Q_\ell$,
\[
 \Pi_{48}=\frac85(2Q-Q^2)=-\frac{\Omega^2}{60}.
\]
The conjugate channel is therefore literally the joined point/line eigenvalue-$2$
carriers, not just a multiplicity ledger.

\paragraph{4328: the irregular Ihara boundary corrected.}
Bass's determinant formula already applies to general finite graphs with $D-I$.  The
Kotani--Sunada theorem bounds non-real poles by
\[
 (d_{\max}-1)^{-1/2}\le |u|\le(d_{\min}-1)^{-1/2},
\]
not by the same expression without square roots, and it does not put every nontrivial
root in that annulus.  For the shipped degree-$2$--$8$ frame graph the reciprocal-root
form is $1\le|\lambda|\le\sqrt7$ for the corresponding non-real roots.  The old
``band filling'' and ``closest irregular Ramanujan'' metric is withdrawn.

\paragraph{4329: a cover certificate retracted by one cycle.}
The stored $\mathbb Z_{2731}$ voltage vector from Pass 4253 has the base eight-cycle
\[
 (3,52,4,48,8,66,18,53)
\]
with signed voltage sum $1328-542+801-1587=0\pmod{2731}$.  The lifted cover therefore
has girth exactly $8$, not at least $16$.  Pass 4260's $4373$-vertex radius-seven tree
ball remains a correct conditional tree count but is not realized by that cover.  The
valid cover ladder is the earlier $\mathbb Z_{359}$ girth-at-least-$14$ cover followed by
the much larger $\mathbb Z_{750019}$ girth-at-least-$18$ cover.

\paragraph{4330--4331: what symmetry and fault detection actually survive.}
Under the literal pair swap
$S(x_0,x_1,x_2,x_3)=(x_2,x_3,x_0,x_1)$, machines A/B/C/D are invariant in the pattern
false/true/false/true.  Machine D therefore removes both the directed time arrow and the
point/frequency structural asymmetry.  Its $0.460435$ peak is symmetric localization,
not residual point/frequency bias.  Affine translations act on the $81$ frames but
descend to neither projective $40$-carrier, so they do not force a point-side load port.

Pass~4331's linear-opcode mismatch census tests incidence after two rails choose among
three selected linear panel opcodes.  It detects $1656/1920=69/80$ differential
single-rail opcode substitutions and $0/960$ shared-control substitutions.  This is
narrower than the Pass~4367/4374 intrinsic flag-register comparator, which detects
$36/39=12/13$ arbitrary one-register substitutions at $q=3$.  The earlier $95.71\%$
number compared each faulty rail with a golden run and was not a dual-rail comparator.

\paragraph{4332--4333: one arithmetic component, and an honest rank.}
The irreducible degree-$33$ factor of the instruction Hashimoto pencil has exactly one
real root in $(5.74,5.75)$ and one in $(3.349,3.350)$ by an exact Sturm certificate.
Thus the global nonbacktracking growth root and the reported slow real mode are Galois
conjugates in one rational primary component.  The rational projector and a
basis-independent localization theorem remain open.

Finally, the complete universal-set census by sizes $4$ through $10$ is
\[
 24,\ 80,\ 114,\ 90,\ 41,\ 10,\ 1,
\]
totalling $360$.  The shipped ISA is not strictly ``12th'': only six sets have lower
$\rho$ at the stated tolerance, while six isomorphic sets share its numerical value and
occupy positions $7$--$12$.  Size is a coarse correlate whose bands overlap, not an
explanation of that tie.

\paragraph{Evidence boundary.}
The three GAP witnesses replay $31/31$, $28/28$, and $17/17$ exact checks.  The panel operators are
three-way relations, not yet a deterministic opcode selector or synthesized B/D RTL.
The exact finite algebra and the retractions above should therefore be read separately
from hardware, continuum, thermodynamic, and phenomenological interpretations.
